% ============================================================
\chapter{LLM Evaluation}
\label{ch:evaluation}
% ============================================================

\begin{intuition}[title=\textbf{Why evaluation is hard}]
Evaluating NLP systems is fundamentally difficult because language is
ambiguous, subjective, and context-dependent. A translation can be
faithful but awkward, a summary can be fluent but unfaithful. Each task
requires tailored metrics, and no automatic metric perfectly captures
human-perceived quality.
\end{intuition}

% ------------------------------------------------------------
\section{Intrinsic vs Extrinsic Evaluation}
\label{sec:eval:types}
% ------------------------------------------------------------

\begin{definition}[Intrinsic Evaluation]
\index{intrinsic evaluation}
\textbf{Intrinsic evaluation} measures the quality of an isolated
component, independently of its downstream application. Examples:
perplexity of a language model, POS tagging accuracy.
\end{definition}

\begin{definition}[Extrinsic Evaluation]
\index{extrinsic evaluation}
\textbf{Extrinsic evaluation} measures the impact of a component on a
complete downstream task. Example: the effect of a better language model
on dialogue system quality.
\end{definition}

\begin{remark}
Intrinsic evaluation is reproducible and inexpensive but does not
guarantee practical utility. Extrinsic evaluation is more relevant but
costlier and harder to isolate.
\end{remark}

% ------------------------------------------------------------
\section{Perplexity}
\label{sec:eval:perplexity}
% ------------------------------------------------------------

\begin{definition}[Perplexity]
\index{perplexity}
The \textbf{perplexity} of a language model $P$ on a test corpus
$w_1, \ldots, w_N$ is:
\[
    \text{PPL} = \exp\left(-\frac{1}{N} \sum_{i=1}^{N}
    \log P(w_i | w_{<i})\right)
    = \exp(H_{\text{cross}})
\]
where $H_{\text{cross}}$ is the cross-entropy between the true
distribution and the model.
\end{definition}

\begin{proposition}[Interpreting Perplexity]
\begin{itemize}
    \item PPL $= \abs{\mathcal{V}}$: the model is uniform (no prediction).
    \item PPL $= 1$: perfect prediction.
    \item In practice, GPT-2 achieves PPL $\approx 22$ on WikiText-103.
\end{itemize}
\end{proposition}

\begin{warning}[title=\textbf{Limitations of Perplexity}]
Perplexity only measures the model's predictive ability on the test
corpus. It does not capture coherence, factual faithfulness, or stylistic
quality. Two models with the same perplexity can generate text of very
different qualities.
\end{warning}

% ------------------------------------------------------------
\section{BLEU}
\label{sec:eval:bleu}
% ------------------------------------------------------------

\begin{definition}[BLEU]
\index{BLEU}
BLEU (Bilingual Evaluation Understudy, Papineni et al., 2002) measures
translation quality by comparing $n$-grams between a candidate $c$ and
reference translations $\{r\}$:
\[
    \text{BLEU} = \text{BP} \cdot \exp\left(\sum_{n=1}^{N}
    \frac{1}{N} \log p_n\right)
\]
where $p_n$ is the modified $n$-gram precision and BP is the brevity
penalty.
\end{definition}

\begin{keyformulas}[title=\textbf{BLEU Components}]
\textbf{Modified $n$-gram precision}:
\[
    p_n = \frac{\sum_{\text{$n$-gram} \in c}
    \min(\text{count}(c), \max_r \text{count}(r))}
    {\sum_{\text{$n$-gram} \in c} \text{count}(c)}
\]
\textbf{Brevity penalty}:
\[
    \text{BP} = \begin{cases}
        1 & \text{if } \abs{c} > \abs{r} \\
        \exp(1 - \abs{r}/\abs{c}) & \text{otherwise}
    \end{cases}
\]
\end{keyformulas}

\begin{remark}
BLEU-4 (with $N=4$) is the standard for machine translation. It
correlates reasonably with human judgment at the corpus level but
poorly at the sentence level.
\end{remark}

% ------------------------------------------------------------
\section{ROUGE}
\label{sec:eval:rouge}
% ------------------------------------------------------------

\begin{definition}[ROUGE]
\index{ROUGE}
ROUGE (Recall-Oriented Understudy for Gisting Evaluation, Lin, 2004)
measures summary quality with a focus on recall:
\begin{itemize}
    \item \textbf{ROUGE-N}: $n$-gram recall:
    $\text{ROUGE-N} = \frac{\sum_{\text{$n$-gram} \in r}
    \min(\text{count}(c), \text{count}(r))}
    {\sum_{\text{$n$-gram} \in r} \text{count}(r)}$
    \item \textbf{ROUGE-L}: longest common subsequence (LCS) based F-measure.
\end{itemize}
\end{definition}

\begin{proposition}[BLEU vs ROUGE]
\begin{center}
\renewcommand{\arraystretch}{1.3}
\begin{tabular}{lcc}
    \toprule
    & \textbf{BLEU} & \textbf{ROUGE} \\
    \midrule
    Orientation & Precision & Recall \\
    Primary use & Translation & Summarization \\
    Granularity & Corpus & Document \\
    \bottomrule
\end{tabular}
\end{center}
\end{proposition}

% ------------------------------------------------------------
\section{BERTScore}
\label{sec:eval:bertscore}
% ------------------------------------------------------------

\begin{definition}[BERTScore]
\index{BERTScore}
BERTScore (Zhang et al., 2020) uses contextual embeddings from BERT
to compute similarity between reference and candidate tokens:
\[
    P_{\text{BERT}} = \frac{1}{\abs{c}} \sum_{c_i \in c}
    \max_{r_j \in r} \mathbf{c}_i^\top \mathbf{r}_j,
    \qquad
    R_{\text{BERT}} = \frac{1}{\abs{r}} \sum_{r_j \in r}
    \max_{c_i \in c} \mathbf{r}_j^\top \mathbf{c}_i
\]
\[
    F_{\text{BERT}} = 2 \cdot \frac{P_{\text{BERT}} \cdot R_{\text{BERT}}}
    {P_{\text{BERT}} + R_{\text{BERT}}}
\]
\end{definition}

\begin{remark}
BERTScore captures semantic similarity rather than exact word matching,
making it more robust to paraphrases. It correlates better with human
judgment than BLEU and ROUGE on many tasks.
\end{remark}

% ------------------------------------------------------------
\section{Human Evaluation}
\label{sec:eval:human}
% ------------------------------------------------------------

\begin{definition}[Human Evaluation Protocols]
\index{human evaluation}
Human evaluation remains the gold standard. Common protocols include:
\begin{enumerate}
    \item \textbf{Likert scale}: annotators rate from 1 to 5 on criteria
    (fluency, coherence, relevance).
    \item \textbf{Pairwise comparison}: ``Which response is better,
    A or B?'' (more reliable than absolute scales).
    \item \textbf{Ranking}: order $n$ outputs from best to worst.
    \item \textbf{Elo rating}: dynamic ranking through iterative
    comparisons (used by Chatbot Arena).
\end{enumerate}
\end{definition}

\begin{definition}[Inter-Annotator Agreement]
\index{inter-annotator agreement}
\textbf{Cohen's kappa} measures agreement between two annotators
corrected for chance:
\[
    \kappa = \frac{P_o - P_e}{1 - P_e}
\]
where $P_o$ is observed agreement and $P_e$ is expected agreement by
chance. $\kappa > 0.8$ indicates excellent agreement.
\end{definition}

\begin{warning}[title=\textbf{Human Evaluation Biases}]
Human evaluation suffers from biases: preference for longer responses,
position bias (first option is often preferred), annotator fatigue, and
cultural variability. Rigorous protocols (randomization, explicit
criteria, calibration) are essential.
\end{warning}

% ------------------------------------------------------------
\section{Benchmarks: GLUE, SuperGLUE, and Beyond}
\label{sec:eval:benchmarks}
% ------------------------------------------------------------

\begin{definition}[GLUE]
\index{GLUE}
The \textbf{GLUE} benchmark (General Language Understanding Evaluation,
Wang et al., 2018) comprises 9 language understanding tasks: textual
entailment (MNLI, RTE, WNLI), similarity (STS-B, MRPC, QQP), sentiment
(SST-2), acceptability (CoLA), and QA (QNLI).
\end{definition}

\begin{definition}[SuperGLUE]
\index{SuperGLUE}
\textbf{SuperGLUE} (Wang et al., 2019) succeeded GLUE with harder tasks:
BoolQ, CB, COPA, MultiRC, ReCoRD, RTE, WiC, WSC. Current models surpass
human performance on SuperGLUE.
\end{definition}

\begin{remark}
Since LLMs have saturated GLUE and SuperGLUE, more challenging
benchmarks have emerged: MMLU (multi-domain knowledge), BIG-Bench
(emergent tasks), HELM (holistic evaluation), and LMSYS Chatbot Arena
(user-based evaluation).
\end{remark}

% ------------------------------------------------------------
\section{LLM-as-a-Judge}
\label{sec:eval:llm_judge}
% ------------------------------------------------------------

\begin{definition}[LLM-as-a-Judge]
\index{LLM-as-a-Judge}
The \textbf{LLM-as-a-Judge} approach uses a powerful language model
(GPT-4, Claude) to evaluate outputs from other models. The judge
receives a structured prompt with evaluation criteria and returns a
score or ranking.
\end{definition}

\begin{attention}[title=\textbf{Leaderboard Pitfalls}]
Leaderboards can be misleading:
\begin{enumerate}
    \item \textbf{Benchmark overfitting}: models are optimized
    specifically for test benchmarks.
    \item \textbf{Data contamination}: test data may appear in training
    corpora.
    \item \textbf{Inappropriate metrics}: a single score masks strengths
    and weaknesses.
    \item \textbf{Lack of diversity}: benchmarks underrepresent certain
    languages and cultures.
\end{enumerate}
\end{attention}

% ------------------------------------------------------------
\section{Python Implementation}
\label{sec:eval:code}
% ------------------------------------------------------------

\begin{codeblock}[title=\textbf{Computing BLEU and ROUGE}]
\begin{minted}{python}
from nltk.translate.bleu_score import sentence_bleu
from rouge_score import rouge_scorer

# BLEU
reference = [["the", "cat", "is", "on", "the", "mat"]]
candidate = ["the", "cat", "is", "sitting", "on", "the", "mat"]
bleu = sentence_bleu(reference, candidate)
print(f"BLEU: {bleu:.4f}")

# ROUGE
scorer = rouge_scorer.RougeScorer(
    ['rouge1', 'rouge2', 'rougeL'], use_stemmer=True)
ref = "the cat is on the mat"
hyp = "the cat is sitting on the mat"
scores = scorer.score(ref, hyp)
for key, val in scores.items():
    print(f"{key}: P={val.precision:.3f} R={val.recall:.3f} "
          f"F={val.fmeasure:.3f}")
\end{minted}
\end{codeblock}

\begin{codeblock}[title=\textbf{BERTScore}]
\begin{minted}{python}
from bert_score import score as bert_score

refs = ["The cat is on the mat."]
hyps = ["The cat is sitting on the mat."]
P, R, F1 = bert_score(hyps, refs, lang="en")
print(f"BERTScore F1: {F1.mean():.4f}")
\end{minted}
\end{codeblock}

% ------------------------------------------------------------
\section{Exercises}
\label{sec:eval:exercises}
% ------------------------------------------------------------

\begin{exercise}[Manual BLEU]
Manually compute the BLEU-2 (bigram) score for the reference ``the small
black cat sleeps'' and the candidate ``the black cat sleeps well''.
Detail the computation of $p_1$, $p_2$, and BP.
\end{exercise}

\begin{exercise}[BERTScore vs BLEU]
On a corpus of 100 (reference, candidate) pairs containing paraphrases,
compare BLEU and BERTScore. Identify cases where the two metrics diverge
and explain why.
\end{exercise}

\begin{exercise}[Human Evaluation Protocol]
Design a human evaluation protocol for a customer service chatbot.
Define criteria, scale, number of annotators, and how to measure
inter-annotator agreement.
\end{exercise}

\begin{exercise}[Benchmark Contamination]
Explain how test data contamination in an LLM's training corpus can
bias results on MMLU. Propose methods to detect and mitigate this
problem.
\end{exercise}
